\chapter{Wprowadzenie}\label{chap:wprowadzenie}

W~ostatnich latach badania w~dziedzinie sztucznej inteligencji przesunęły się
z~analizy pojedynczych modeli na złożone systemy. Duże modele językowe
(ang.~\textit{Large Language Models}, LLM) oparte na architekturze Transformer
osiągają wysoką skuteczność w~generowaniu tekstu na podstawie ogólnej
wiedzy~\cite{vaswani2017attention}, jednakże w~złożonych zadaniach wymagających
wieloetapowego wnioskowania i~wiedzy dziedzinowej generacja tekstu jest
niewystarczająca~\cite{brown2020language}. Rozwiązanie takich zadań wymaga
wyszukiwania i~integracji informacji oraz formułowania wniosków na podstawie
pobranych źródeł. Zmiana paradygmatu ukształtowała dziedzinę architektur
agentowych, wykorzystywania zewnętrznych narzędzi~\cite{schick2023toolformer},
generowania wspomaganego wyszukiwaniem (ang.~\textit{Retrieval-Augmented
Generation}, RAG)~\cite{lewis2020retrieval} oraz protokołów
orkiestracji~\cite{wang2024survey}. W~literaturze brakuje jednak kontrolowanych
porównań wskazujących, kiedy dana architektura orkiestracji poprawia jakość
wyników, a~kiedy generuje jedynie dodatkowy narzut obliczeniowy i~koszty
operacyjne.

\section{Motywacja i~problem badawczy}\label{sec:motywacja-i-problem-badawczy}

Projektowanie systemów agentowych opartych na dużych modelach językowych do
odpowiadania na pytania na podstawie obszernych zbiorów dokumentów pozostaje
otwartym problemem badawczym. Modele wykazują wyższą skuteczność po wyposażeniu
w~odpowiednie narzędzia i~ustrukturyzowany schemat
wnioskowania~\cite{wang2024survey}.

W~architekturze ReAct zaobserwowano, że przeplatanie etapów wnioskowania
(ang.~\textit{reasoning traces}) z~wywołaniami narzędzi (ang.~\textit{tool
calls}) zmniejsza liczbę halucynacji i~zwiększa skuteczność w~porównaniu do
podejścia łańcucha myśli (ang.~\textit{chain-of-thought})~\cite{yao2022react}.
Systemy wieloagentowe, takie jak AutoGen, udowodniły skuteczność podziału
zadań~\cite{wu2024autogen}, natomiast system MetaGPT potwierdził wpływ
specjalizacji ról na wyniki w~złożonych zadaniach~\cite{hong2023metagpt}.

Architektury te rzadko podlegają jednak bezpośrednim porównaniom. Nowe
rozwiązania zazwyczaj zestawia się z~modelem bazowym -- wyposażonym w~prostą
instrukcję systemową (ang.~\textit{system prompt}) lub w~architekturze
jednoprzebiegowego systemu RAG (ang.~\textit{single-pass RAG}). Podejście to
uniemożliwia obiektywną ocenę struktur agentowych w~identycznych warunkach --
przy wykorzystaniu tego samego modelu bazowego, zestawu narzędzi, korpusu
danych i~metryk ewaluacyjnych.

Brak ujednoliconych testów utrudnia syntezę wiedzy naukowej oraz projektowanie
systemów produkcyjnych. Według Wang i~in. porównawcza ewaluacja topologii
agentów w~stałych warunkach pozostaje problemem otwartym~\cite{wang2024survey}.
Xi i~in. wskazują, że niespójność zadań, modeli bazowych i~metryk uniemożliwia
porównywanie wyników pomiędzy publikacjami~\cite{xi2025rise}.

Niniejsza praca podejmuje problem braku kontrolowanego porównania izolującego
wpływ struktury agenta na poprawność odpowiedzi oraz koszty operacyjne podczas
analizy ustrukturyzowanych dokumentów regulacyjnych. Badaniu poddano zależność
między złożonością systemów wieloagentowych a~jakością uzyskiwanych odpowiedzi
w~kontrolowanym środowisku testowym.

Jako środowisko testowe wybrano polską dokumentację ubezpieczeniową. Ogólne
Warunki Ubezpieczenia (OWU) oraz dokumenty zawierające informacje o~produkcie
ubezpieczeniowym (IPID) stanowią trudną klasę tekstów dla systemów pytań
i~odpowiedzi (ang.~\textit{Question Answering}, QA). Dokumenty te
charakteryzują się specjalistycznym językiem prawniczo-ubezpieczeniowym,
zawierają rozbudowane struktury warunków i~wyłączeń, a~kluczowe informacje są
w~nich rozproszone. Wymusza to stosowanie mechanizmów wyszukiwania
wieloskokowego i~wnioskowania warunkowego, co pozwala na miarodajne
zróżnicowanie badanych architektur agentowych.

Rysunek~\ref{fig:projekt-eksperymentu} przedstawia główne założenie
eksperymentu: jedyną zmienną niezależną jest topologia agenta, podczas gdy
model, narzędzia, korpus i~warstwa ewaluacyjna pozostają wspólne dla wszystkich
wariantów.

\begin{figure}[H]
  \caption{Projekt eksperymentu}\label{fig:projekt-eksperymentu}
  \centering
  \begin{adjustbox}{max width=\textwidth}
    \begin{tikzpicture}[
        node distance=0.4cm and 0.25cm,
        archbox/.style={
          rectangle, draw, rounded corners=3pt,
          text width=2.0cm, minimum height=1.25cm,
          align=center, font=\small,
          fill=violet!8, draw=violet!50
        },
        grouplabel/.style={
          font=\small\bfseries,
          fill=white, inner xsep=6pt, inner ysep=2pt
        },
        groupbox/.style={
          rectangle, draw=gray!50, dashed,
          rounded corners=5pt, inner sep=8pt
        },
        midbox/.style={
          rectangle, draw, rounded corners=4pt,
          inner sep=9pt, align=center, font=\small
        },
        arrow/.style={->, >=stealth, draw=gray!60}
      ]

      \node[archbox] (arch1) {Determi-\\nistyczna};
      \node[archbox, right=of arch1] (arch2) {Planer--\\wykonawca};
      \node[archbox, right=of arch2] (arch3) {Router--\\specjalista};
      \node[archbox, right=of arch3] (arch4) {Tablicowa};
      \node[archbox, right=of arch4] (arch5) {Hierar-\\chiczna};
      \node[archbox, right=of arch5] (arch6) {ReAct};

      \node[groupbox, fit=(arch1)(arch6)] (var_group) {};
      \node[grouplabel] at (var_group.north)
      {Zmienna niezależna: topologia systemu};

      \node[midbox, fill=gray!7, draw=gray!40,
      below=0.9cm of var_group] (baseline) {%
        \textbf{Punkt odniesienia}\\[3pt]
        Stały model bazowy \(\cdot\)
        Stały zestaw narzędzi\\[2pt]
        Stały indeks wektorowy\(\cdot\)
        Stały zestaw testowy pytań
      };

      \node[midbox, fill=teal!7, draw=teal!40,
      below=0.9cm of baseline] (eval) {%
        \textbf{Warstwa ewaluacji}\\[3pt]
        Wierność \(\cdot\)
        Poprawność \(\cdot\)
        Zwięzłość  \(\cdot\)
        Trafność pobranych dokumentu\\[2pt]
        Opóźnienie \(\cdot\)
        Zużycie tokenów
      };

      \draw[arrow] (var_group.south) -- (baseline.north);
      \draw[arrow] (baseline.south) -- (eval.north);

    \end{tikzpicture}
  \end{adjustbox}
  \zrodlo{opracowanie własne}
\end{figure}

\clearpage{}

\section{Pytania badawcze}\label{sec:pytania-badawcze}

Pytania badawcze wynikają bezpośrednio z~luki opisanej
w~\hyperref[sec:motywacja-i-problem-badawczy]{sekcji~\ref*{sec:motywacja-i-problem-badawczy}}
i~odpowiadają czterem obszarom porównania: jakości odpowiedzi, stosunkowi
kosztu do jakości, wpływowi strategii rozwiązywania pytań trudnych oraz
charakterystyce błędów.

\begin{description}
  \item[Pytanie badawcze 1 (PB1):] W~jakim stopniu poszczególne architektury różnią się
    pod względem wierności źródeł? Weryfikacji poddano poprawność merytoryczną
    i~trafność pobieranych dokumentów w~kontrolowanym środowisku ze stałym modelem,
    zestawem narzędzi, korpusem i~metrykami.

  \item[Pytanie badawcze 2 (PB2):] Która topologia agentowa wykazuje optymalny stosunek
    jakości do kosztów operacyjnych? Koszty operacyjne wyrażono poprzez opóźnienie
    oraz zużycie tokenów. Analizie poddano również zmiany tego stosunku
    w~zależności od stopnia trudności zadanego pytania.

  \item[Pytanie badawcze 3 (PB3):] Jak mechanizmy planowania, trasowania
    (ang.~\textit{routing}) i~podziału zadań wpływają na rozwiązywanie złożonych
    problemów, w~tym na zestawianie warunków ochrony z~wyjątkami w~dokumentach?
    Skuteczność mechanizmów porównano z~układami realizującymi wnioskowanie
    w~pojedynczej ścieżce.

  \item[Pytanie badawcze 4 (PB4):] Jakie rodzaje błędów wykazują poszczególne struktury
    podczas analizy tekstów ustrukturyzowanych? Określono potencjalne źródła tych
    błędów, w~tym ograniczenia okna kontekstowego, zakłócenia komunikacji między
    modułami oraz kaskadową propagację błędów.
\end{description}

Tabela~\ref{tab:zestawienie-pytan-badawczych-hipotez-i-zrodel-danych} zestawia
każde pytanie badawcze z~odpowiadającą mu hipotezą oraz wskazuje sekcje,
w~których zgromadzono dane niezbędne do jej weryfikacji.

\begin{table}[H]
  \caption{Zestawienie pytań badawczych, hipotez i źródeł danych}\label{tab:zestawienie-pytan-badawczych-hipotez-i-zrodel-danych}
  \centering
  \small
  \renewcommand{\arraystretch}{1.3}
  \begin{tabularx}{\textwidth}{@{} l X X @{}}
    \toprule
    \textbf{Pytanie} & \textbf{Hipoteza}                                                         & \textbf{Źródło dowodów} \\
    \midrule
    PB1
    & Bardziej złożone topologie zapewniają wyższą wierność odpowiedzi
    w~porównaniu do prostych układów jednoagentowych
    & Oceny jakościowe i~ilościowe, przekroje według trudności pytań
    --- \hyperref[sec:wyniki-jakosciowe]{sekcja~\ref*{sec:wyniki-jakosciowe}}                                              \\
    PB2
    & Wzrost kosztów operacyjnych w~systemach wieloagentowych nie przekłada się
    proporcjonalnie na przyrost jakości odpowiedzi
    & Analiza zużycia tokenów i~opóźnień, oceny jakościowe
    --- \hyperref[sec:wyniki-operacyjne]{sekcja~\ref*{sec:wyniki-operacyjne}},
    \hyperref[sec:wyniki-jakosciowe]{\ref*{sec:wyniki-jakosciowe}}                                                         \\
    PB3
    & Dekompozycja zadań w~różnym stopniu poprawia skuteczność poszczególnych
    architektur w~pytaniach wieloskokowych
    & Przekroje według trudności pytań, analiza jakościowa przebiegów
    --- \hyperref[sec:wyniki-jakosciowe]{sekcja~\ref*{sec:wyniki-jakosciowe}}                                              \\
    PB4
    & Różne architektury agentowe wykazują odmienne charakterystyki błędów
    przy identycznych zadaniach
    & Analiza jakościowa przebiegów, klasyfikacja niepowodzeń
    --- \hyperref[subsec:charakterystyka-niepowodzen]{sekcja~\ref*{subsec:charakterystyka-niepowodzen}}                    \\
    \bottomrule
  \end{tabularx}
  \zrodlo{opracowanie własne}
\end{table}

\section{Cele badawcze}\label{sec:cele-badawcze}

Odpowiedź na postawione pytania badawcze wymagała zrealizowania czterech
szczegółowych celów.

\begin{description}
  \item[Cel badawczy 1 (CB1):] Opracowano środowisko orkiestracji oparte na frameworku
    LangGraph i~platformie ewaluacyjnej Arize Phoenix, w~którym zaimplementowano
    sześć architektur agentowych. Umożliwiło to izolację topologii systemu jako
    jedynej zmiennej niezależnej w~badaniu.

  \item[Cel badawczy 2 (CB2):] Skonstruowano zestaw testowy dla branży
    ubezpieczeniowej, składający się z~90 pytań opartych na dokumentach OWU i~IPID,
    odzwierciedlających rzeczywiste zapytania użytkowników.

  \item[Cel badawczy 3 (CB3):] Przeprowadzono ewaluację wszystkich wariantów przy
    użyciu jednolitych metryk jakościowych (wierność, poprawność) i~operacyjnych.

  \item[Cel badawczy 4 (CB4):] Przeanalizowano uzyskane dane w~celu identyfikacji cech
    charakterystycznych badanych struktur oraz sformułowania rekomendacji dla
    projektantów systemów produkcyjnych.
\end{description}

\section{Zakres i~ograniczenia}\label{sec:zakres}

Eksperyment przeprowadzono w~zamkniętym środowisku testowym. Bazę wiedzy
stanowiły 22~polskojęzyczne dokumenty ubezpieczeniowe, podzielone na 5108
fragmentów tekstowych (ang.~\textit{chunks}). Architektury generowały
odpowiedzi wyłącznie na podstawie tej bazy, bez dostępu do internetu
i~zewnętrznych źródeł, co zagwarantowało izolację topologii agenta jako
zmiennej niezależnej. Procedura testowa objęła zestaw 90~pytań oraz sześć
architektur agentowych. Wszystkie systemy korzystały z~tego samego modelu
językowego, identycznych narzędzi oraz identycznej bazy wektorowej,
a~instrukcje systemowe zdefiniowano w~języku angielskim.

Wyniki eksperymentu poddano analizie z~trzech perspektyw. Pierwsza obejmuje
pomiary ilościowe (wierność, poprawność, zwięzłość, trafność pobieranych
dokumentów oraz koszty operacyjne) dla wszystkich architektur. Druga
perspektywa dotyczy analizy jakościowej przebiegów: sposobu konstrukcji
odpowiedzi, utraty spójności kontekstu oraz charakterystycznych wzorców błędów.
Trzecia perspektywa obejmuje rekomendacje praktyczne dotyczące doboru
architektury, progu opłacalności orkiestracji oraz roli przygotowania danych
w~zapewnieniu skuteczności systemu.

W~badaniu oddzielono celowe decyzje projektowe od obiektywnych ograniczeń.
Wybór jednego modelu bazowego i~stałego zestawu narzędzi stanowił warunek
konieczny do izolacji wpływu struktury agentowej. Z~tego powodu zrezygnowano
z~dostrajania modeli (ang.~\textit{fine-tuning}), obsługi wielu języków,
indeksowania w~czasie rzeczywistym oraz testów z~udziałem użytkowników
końcowych.

Badanie ma trzy główne ograniczenia. Po pierwsze, zestaw 90~pytań nie
wyczerpuje pełnego zakresu zapytań w~środowisku produkcyjnym. Po drugie, każdą
architekturę uruchomiono jednokrotnie dla każdego pytania. Wszystkie wywołania
(540 przebiegów agentowych i~ewaluacja metodą \textit{LLM-as-a-judge})
zrealizowano przez płatne API OpenAI, co przy ograniczonym budżecie wykluczyło
powtórzenia. Użycie modeli open-source lokalnie odrzucono ze względu na
wymagania sprzętowe (infrastruktura GPU) dla dużych modeli oraz
niewystarczającą jakość analizy polskich tekstów prawniczych przez modele
mniejsze. Ustawienie temperatury modelu na zero oraz powtarzalny potok
wyszukiwania miały na celu minimalizację wpływu losowości. Trzecim
ograniczeniem jest automatyczna ewaluacja z~użyciem modelu jako sędziego
(ang.~\textit{LLM-as-a-judge}), co wprowadza ryzyko błędów poznawczych modelu
i~nie zastępuje weryfikacji eksperckiej. Wyników nie należy uogólniać na inne
dziedziny bez dodatkowych badań.

\section{Struktura pracy}\label{sec:struktura-pracy}

Układ rozdziałów odpowiada kolejnym fazom badań.

\begin{description}
  \item[\Cref{chap:podstawy} -- Podstawy teoretyczne:] Omawia kluczowe pojęcia z~zakresu architektur agentowych i~systemów generowania wspomaganego wyszukiwaniem oraz charakteryzuje polską dokumentację ubezpieczeniową jako dziedzinę badawczą.

  \item[\Cref{chap:metodologia} -- Metodologia:] Przedstawia założenia eksperymentu, sposób budowy korpusu i~zestawu testowego, przyjęte kryteria i~metryki oceny oraz warunki zapewniające porównywalność wyników między architekturami.

  \item[\Cref{chap:implementacja} -- Implementacja:] Opisuje sposób realizacji sześciu architektur agentowych, zastosowane narzędzia, budowę bazy wektorowej oraz protokół ewaluacji przyjęty na potrzeby eksperymentu.

  \item[\Cref{chap:eksperyment-i-wyniki} -- Eksperyment i wyniki:] Prezentuje wyniki 540~uruchomień z~podziałem na poziomy trudności pytań wraz z~analizą kosztów operacyjnych i~zestawieniem zaobserwowanych zachowań architektur.

  \item[\Cref{chap:wnioski-i-zakonczenie} -- Wnioski i zakończenie:] Odpowiada na postawione pytania badawcze, formułuje rekomendacje dla praktyków projektujących systemy agentowe oraz wskazuje kierunki dalszych badań.
\end{description}
